\documentclass[11pt]{article}
\usepackage{fontspec}
\setmainfont{Latin Modern Roman}
\usepackage{amsmath,amssymb,mathtools,unicode-math}
\setmathfont{Latin Modern Math}
\usepackage{xurl,hyperref}
\hypersetup{colorlinks=true,urlcolor=blue,pdftitle={Problem statement}}
\newcommand{\sref}[2]{\hyperlink{source:#1:#2}{[S#2]}}
\newcommand{\cataloguescope}[1]{\paragraph{Mathematical question.}#1}
\begin{document}
\section{Algebraic Stories — Generic $k$-rank}\label{statement}

\subsection{Definitions and mathematical statement}

\paragraph{Shared notation.}$\mathbb N=\{1,2,\ldots\}$, and $\mathbb Z,\mathbb Q,\mathbb R,\mathbb C$ denote the integers, rationals, reals and complex numbers. Any other conventions are specified locally.


Work over $\mathbb C$. For $n\geq1$, let $S=\mathbb C[x_1,\ldots,x_n]$, with $S_e$ its vector space of homogeneous polynomials of total degree $e$. For $k,d\geq1$ and $f\in S_{kd}$, define
\[\operatorname{rk}_k(f)=\min\left\{s\geq0:f=\sum_{j=1}^{s}g_j^k,\quad g_j\in S_d\right\}.
\]
A general form means one whose coefficient vector lies in a suitable nonempty Zariski open subset of $S_{kd}$. Write $\operatorname{rk}_k^\circ(kd,n)$ for the value on general forms.
\paragraph{Conjecture with its original parameter restriction.}
For all integers $k,n,d\geq2$, is
\[
\operatorname{rk}_k^\circ(kd,n)=
\begin{cases}
\min\{s\in\mathbb N:s\binom{n+d-1}{n-1}-\binom{s}{2}\geq\binom{n+2d-1}{n-1}\},&k=2,\\
\min\{s\in\mathbb N:s\binom{n+d-1}{n-1}\geq\binom{n+kd-1}{n-1}\},&k\geq3?
\end{cases}
\] \sref{statement}{2}\sref{statement}{3}
\subsection{Short English statement}
For powers of forms of degree at least two, test the proposed generic-rank formula, including the correction for sums of squares.
\subsection{Source scope and status}
The restriction $k,n,d\geq2$ is explicit in Lundqvist--Oneto--Reznick--Shapiro, Conjecture 1.2. The later \emph{Algebraic Stories} display and the previous catalogue wording omitted the exclusion of $d=1$. \sref{statement}{2}\sref{statement}{3}

The unrestricted catalogue variant is false: at $(k,d,n)=(4,1,3)$ its formula gives $5$, whereas a general complex ternary quartic has fourth-power rank $6$. This is a classical Alexander--Hirschowitz exception, already recorded in \emph{Algebraic Stories}, Theorem 1.4. The reviewed contribution supplies an explicit proof of this example. It corrects the catalogue's scope and does not disprove or resolve the $d\geq2$ conjecture stated above. \sref{statement}{2}\sref{statement}{4}
\subsection{Sources}
\begin{description}
\item[\hypertarget{source:statement:1}{S1}] ULAM.AI and the UnsolvedMath Contributors, UnsolvedMath: A Curated Collection of Open Mathematics Problems, catalogue entry. CC BY 4.0. \url{https://huggingface.co/datasets/ulamai/UnsolvedMath}.
\item[\hypertarget{source:statement:2}{S2}] R. Fröberg, S. Lundqvist, A. Oneto and B. Shapiro, \emph{Algebraic Stories from One and from the Other Pockets}, Arnold Mathematical Journal 4 (2018), 137--160, §1.1, Theorem 1.4 and Conjecture 1.6, pp.~139--140. \url{https://doi.org/10.1007/s40598-018-0088-z}.
\item[\hypertarget{source:statement:3}{S3}] S. Lundqvist, A. Oneto, B. Reznick and B. Shapiro, \emph{On generic and maximal k-ranks of binary forms}, arXiv:1711.05014v1 (2017), Conjecture 1.2, p.~2. \url{https://arxiv.org/abs/1711.05014v1}.
\item[\hypertarget{source:statement:4}{S4}] We-Solve research, contribution \#30, \emph{\texttt{prb\_c7dda75cdc814d0e8538a125ab32935f}: a classical counterexample to the catalogue's unrestricted formula}; catalogue correction and explicit ternary-quartic proof, reviewed 18 October 2026. \url{https://github.com/We-Solve/research/issues/30}.
\end{description}
\label{end:statement}
\end{document}
